\question
% Тема: Минимизация булевых функций
Комендант общежития ИТМО Галина Петровна устала объяснять новому вахтёру Гоше, кого пускать в общежитие, и решила заказать «умный турникет». Правила у неё такие: жильца с действующим пропуском днём пускаем всегда; ночью (с 23:00 до 6:00) жильца с пропуском пускаем, только если у него вдобавок есть ночное разрешение; гостя без пропуска пускаем только днём и только если его сопровождает староста этажа.

Введём переменные: $P$ — у пришедшего действующий пропуск, $S$ — его сопровождает староста, $N$ — сейчас ночь, $R$ — есть ночное разрешение.
1. Запишите булеву функцию $F(P, S, N, R)$ «турникет открывается» прямо по правилам Галины Петровны.
2. Постройте таблицу истинности и выпишите СДНФ. Сколько в ней минтермов?
3. Гоше нужно объяснение попроще: найдите минимальную ДНФ и сформулируйте её словами.

---------------

Автор -- Ковалев Евгений